\section{Mode Operasi Block Cipher}
\label{sec:mode-operasi}

Cipher blok hanya mampu mengenkripsi satu blok data pada satu waktu, sedangkan
pesan yang hendak dikirim hampir selalu lebih panjang daripada satu blok.
Bagaimana blok-blok itu dirangkai, apakah setiap blok berdiri sendiri atau
bergantung pada blok sebelumnya, dan bagaimana blok terakhir yang tidak lengkap
ditangani --- semuanya diatur oleh \textbf{mode operasi}. Mode operasi bukan
bagian dari algoritma cipher blok itu sendiri; AES tetap AES, tetapi AES-ECB
dan AES-CBC berperilaku sangat berbeda terhadap keamanan dan terhadap pola data.

Ada lima mode operasi klasik yang dibakukan dan dibahas di sini: \textit{Electronic
Code Book} (ECB), \textit{Cipher Block Chaining} (CBC), \textit{Cipher Feedback}
(CFB), \textit{Output Feedback} (OFB), dan \textit{Counter Mode} (CTR). Kelimanya
memakai fungsi enkripsi $E_K$ yang sama; yang berbeda adalah cara masukan bagi
$E_K$ dibentuk dan cara keluarannya dipakai.

\subsection{Electronic Code Book (ECB)}

Mode ECB adalah mode yang paling langsung. Setiap blok plainteks $P_i$
dienkripsi secara individual dan independen dari blok lainnya:
\begin{equation}
C_i = E_K(P_i) \qquad \text{dan} \qquad P_i = D_K(C_i).
\label{eq:ecb}
\end{equation}
Tidak ada keadaan yang dibawa dari satu blok ke blok berikutnya. Gambar~\ref{fig:mode-ecb-cbc}
membandingkan susunan ECB dengan CBC yang dibahas setelah ini.

% Perbandingan mode ECB dan CBC.
% Dirujuk dari section-02.tex dengan % Perbandingan mode ECB dan CBC.
% Dirujuk dari section-02.tex dengan % Perbandingan mode ECB dan CBC.
% Dirujuk dari section-02.tex dengan \input{figures/mode-ecb-cbc}.
\begin{figure}[htbp]
    \centering
    \begin{tikzpicture}[
        kotak/.style={draw, rounded corners, align=center, font=\small, inner sep=3pt,
                      minimum width=0.95cm, minimum height=0.6cm},
        blok/.style={kotak, fill=blue!6},
        enc/.style={kotak, fill=orange!12, minimum width=1.1cm},
        cip/.style={kotak, fill=green!10},
        xor/.style={draw, circle, inner sep=1pt, font=\small, fill=yellow!18,
                    minimum size=0.5cm},
        panah/.style={-{Latex[length=2mm]}, thick},
    ]
        % ================= ECB =================
        \node[font=\small\bfseries] at (3,0.9) {Electronic Code Book (ECB)};
        \foreach \i/\x in {1/0, 2/3, 3/6} {
            \node[blok] (p\i) at (\x,0)   {$P_{\i}$};
            \node[enc]  (e\i) at (\x,-1.3) {$E_K$};
            \node[cip]  (c\i) at (\x,-2.6) {$C_{\i}$};
            \draw[panah] (p\i) -- (e\i);
            \draw[panah] (e\i) -- (c\i);
        }
        \node[font=\scriptsize, align=center] at (3,-3.5)
            {setiap blok dienkripsi sendiri-sendiri: blok plainteks yang sama\\ selalu menghasilkan cipherteks yang sama};

        % ================= CBC =================
        \node[font=\small\bfseries] at (3,-5.0) {Cipher Block Chaining (CBC)};
        \foreach \i/\x in {1/0, 2/3, 3/6} {
            \node[blok] (q\i) at (\x,-6.0) {$P_{\i}$};
            \node[xor]  (x\i) at (\x,-7.0) {$\oplus$};
            \node[enc]  (f\i) at (\x,-8.0) {$E_K$};
            \node[cip]  (d\i) at (\x,-9.0) {$C_{\i}$};
            \draw[panah] (q\i) -- (x\i);
            \draw[panah] (x\i) -- (f\i);
            \draw[panah] (f\i) -- (d\i);
        }
        \node[blok, fill=gray!12] (iv) at (-2.6,-7.0) {$IV$};
        \draw[panah] (iv) -- (x1);
        \draw[panah, dashed] (d1.east) -- node[sloped, above, font=\scriptsize] {$C_1$} (x2.west);
        \draw[panah, dashed] (d2.east) -- node[sloped, above, font=\scriptsize] {$C_2$} (x3.west);
        \node[font=\scriptsize, align=center] at (3,-9.9)
            {setiap blok di-XOR dengan cipherteks sebelumnya, sehingga blok yang sama\\ menghasilkan cipherteks berbeda sesuai posisinya};
    \end{tikzpicture}
    \caption{Perbandingan mode ECB yang independen dengan mode CBC yang merantai blok melalui operasi XOR}
    \label{fig:mode-ecb-cbc}
\end{figure}
.
\begin{figure}[htbp]
    \centering
    \begin{tikzpicture}[
        kotak/.style={draw, rounded corners, align=center, font=\small, inner sep=3pt,
                      minimum width=0.95cm, minimum height=0.6cm},
        blok/.style={kotak, fill=blue!6},
        enc/.style={kotak, fill=orange!12, minimum width=1.1cm},
        cip/.style={kotak, fill=green!10},
        xor/.style={draw, circle, inner sep=1pt, font=\small, fill=yellow!18,
                    minimum size=0.5cm},
        panah/.style={-{Latex[length=2mm]}, thick},
    ]
        % ================= ECB =================
        \node[font=\small\bfseries] at (3,0.9) {Electronic Code Book (ECB)};
        \foreach \i/\x in {1/0, 2/3, 3/6} {
            \node[blok] (p\i) at (\x,0)   {$P_{\i}$};
            \node[enc]  (e\i) at (\x,-1.3) {$E_K$};
            \node[cip]  (c\i) at (\x,-2.6) {$C_{\i}$};
            \draw[panah] (p\i) -- (e\i);
            \draw[panah] (e\i) -- (c\i);
        }
        \node[font=\scriptsize, align=center] at (3,-3.5)
            {setiap blok dienkripsi sendiri-sendiri: blok plainteks yang sama\\ selalu menghasilkan cipherteks yang sama};

        % ================= CBC =================
        \node[font=\small\bfseries] at (3,-5.0) {Cipher Block Chaining (CBC)};
        \foreach \i/\x in {1/0, 2/3, 3/6} {
            \node[blok] (q\i) at (\x,-6.0) {$P_{\i}$};
            \node[xor]  (x\i) at (\x,-7.0) {$\oplus$};
            \node[enc]  (f\i) at (\x,-8.0) {$E_K$};
            \node[cip]  (d\i) at (\x,-9.0) {$C_{\i}$};
            \draw[panah] (q\i) -- (x\i);
            \draw[panah] (x\i) -- (f\i);
            \draw[panah] (f\i) -- (d\i);
        }
        \node[blok, fill=gray!12] (iv) at (-2.6,-7.0) {$IV$};
        \draw[panah] (iv) -- (x1);
        \draw[panah, dashed] (d1.east) -- node[sloped, above, font=\scriptsize] {$C_1$} (x2.west);
        \draw[panah, dashed] (d2.east) -- node[sloped, above, font=\scriptsize] {$C_2$} (x3.west);
        \node[font=\scriptsize, align=center] at (3,-9.9)
            {setiap blok di-XOR dengan cipherteks sebelumnya, sehingga blok yang sama\\ menghasilkan cipherteks berbeda sesuai posisinya};
    \end{tikzpicture}
    \caption{Perbandingan mode ECB yang independen dengan mode CBC yang merantai blok melalui operasi XOR}
    \label{fig:mode-ecb-cbc}
\end{figure}
.
\begin{figure}[htbp]
    \centering
    \begin{tikzpicture}[
        kotak/.style={draw, rounded corners, align=center, font=\small, inner sep=3pt,
                      minimum width=0.95cm, minimum height=0.6cm},
        blok/.style={kotak, fill=blue!6},
        enc/.style={kotak, fill=orange!12, minimum width=1.1cm},
        cip/.style={kotak, fill=green!10},
        xor/.style={draw, circle, inner sep=1pt, font=\small, fill=yellow!18,
                    minimum size=0.5cm},
        panah/.style={-{Latex[length=2mm]}, thick},
    ]
        % ================= ECB =================
        \node[font=\small\bfseries] at (3,0.9) {Electronic Code Book (ECB)};
        \foreach \i/\x in {1/0, 2/3, 3/6} {
            \node[blok] (p\i) at (\x,0)   {$P_{\i}$};
            \node[enc]  (e\i) at (\x,-1.3) {$E_K$};
            \node[cip]  (c\i) at (\x,-2.6) {$C_{\i}$};
            \draw[panah] (p\i) -- (e\i);
            \draw[panah] (e\i) -- (c\i);
        }
        \node[font=\scriptsize, align=center] at (3,-3.5)
            {setiap blok dienkripsi sendiri-sendiri: blok plainteks yang sama\\ selalu menghasilkan cipherteks yang sama};

        % ================= CBC =================
        \node[font=\small\bfseries] at (3,-5.0) {Cipher Block Chaining (CBC)};
        \foreach \i/\x in {1/0, 2/3, 3/6} {
            \node[blok] (q\i) at (\x,-6.0) {$P_{\i}$};
            \node[xor]  (x\i) at (\x,-7.0) {$\oplus$};
            \node[enc]  (f\i) at (\x,-8.0) {$E_K$};
            \node[cip]  (d\i) at (\x,-9.0) {$C_{\i}$};
            \draw[panah] (q\i) -- (x\i);
            \draw[panah] (x\i) -- (f\i);
            \draw[panah] (f\i) -- (d\i);
        }
        \node[blok, fill=gray!12] (iv) at (-2.6,-7.0) {$IV$};
        \draw[panah] (iv) -- (x1);
        \draw[panah, dashed] (d1.east) -- node[sloped, above, font=\scriptsize] {$C_1$} (x2.west);
        \draw[panah, dashed] (d2.east) -- node[sloped, above, font=\scriptsize] {$C_2$} (x3.west);
        \node[font=\scriptsize, align=center] at (3,-9.9)
            {setiap blok di-XOR dengan cipherteks sebelumnya, sehingga blok yang sama\\ menghasilkan cipherteks berbeda sesuai posisinya};
    \end{tikzpicture}
    \caption{Perbandingan mode ECB yang independen dengan mode CBC yang merantai blok melalui operasi XOR}
    \label{fig:mode-ecb-cbc}
\end{figure}


Nama \textit{code book} berasal dari sifatnya yang menyerupai buku kode:
untuk setiap kunci $K$ yang tetap, pemetaan $P \mapsto E_K(P)$ adalah tabel
yang pasti. Bila ukuran blok $n$ bit, maka tabel itu berisi $2^n$ entri, dan
untuk setiap kunci yang berbeda terdapat tabel yang berbeda pula, sehingga
seluruhnya ada $2^m$ tabel untuk kunci $m$ bit. Menyimpan tabel semacam itu
jelas tidak mungkin: untuk blok 64 bit saja, satu tabel berisi $2^{64}$ entri.
Karena itu ECB tidak pernah diimplementasikan sebagai tabel, melainkan sebagai
pemanggilan fungsi $E_K$ --- tetapi sifat matematisnya tetap sama.

Kelebihan ECB bersumber dari sifat independennya. Karena setiap blok berdiri
sendiri, blok-blok tidak perlu diproses berurutan: blok pertama, blok terakhir,
lalu blok di tengah dapat dienkripsi dalam urutan apa pun, dan setiap blok dapat
diakses secara acak tanpa harus mendekripsi blok-blok sebelumnya. Sifat ini
berguna untuk mengenkripsi arsip atau rekaman basis data yang dibaca secara
acak. Kelebihan kedua adalah pelokalan galat: kesalahan satu bit atau lebih pada
sebuah blok cipherteks hanya merusak blok plainteks yang bersangkutan, dan
blok-blok lain tidak terpengaruh.

Kelemahan ECB justru merupakan cerminan dari kelebihannya, dan kelemahan itu
serius. Karena blok plainteks yang sama selalu menghasilkan blok cipherteks yang
sama, pola yang berulang pada plainteks tetap tampak pada cipherteks. Pada teks,
bagian yang berulang --- misalnya deretan spasi yang panjang --- akan muncul
sebagai deretan blok cipherteks yang identik pula, sehingga cipherteks dapat
diserang dengan analisis frekuensi tanpa perlu mengetahui kuncinya.

Kelemahan kedua lebih halus tetapi sama berbahayanya: pihak lawan dapat
memanipulasi cipherteks untuk menipu penerima. Andaikan seseorang mengirim pesan
\begin{center}
``Uang ditransfer \textbf{lima} satu juta rupiah''
\end{center}
dan kriptanalis mengetahui bahwa ukuran blok adalah 2 karakter serta mengetahui
posisi blok yang memuat kata ``lima''. Tanpa mengetahui kuncinya sama sekali,
kriptanalis dapat membuang kedua blok cipherteks yang bersesuaian dengan kata itu,
sehingga cipherteks yang diterima penerima tinggal berbunyi
``Uang ditransfer satu juta rupiah''. Penerima mendekripsinya dengan kunci yang
benar dan memperoleh pesan yang bermakna, lalu menyimpulkan bahwa uang yang
dikirim kepadanya hanya satu juta rupiah. Perhatikan bahwa tidak ada kunci yang
dipecahkan di sini: yang dieksploitasi adalah sifat ECB bahwa tiap blok dapat
dibuang atau disusun ulang tanpa merusak blok lainnya.

Jalan keluar dari kelemahan itu adalah membuat enkripsi setiap blok bergantung
pada seluruh blok sebelumnya, sehingga memindahkan atau membuang satu blok
merusak keseluruhan pesan. Prinsip itulah yang mendasari mode CBC.

\subsection{Cipher Block Chaining (CBC)}

Mode CBC mengikat blok-blok plainteks menjadi satu rantai. Sebelum dienkripsi,
setiap blok plainteks di-XOR terlebih dahulu dengan blok cipherteks sebelumnya:
\begin{equation}
C_0 = IV, \qquad C_i = E_K(P_i \oplus C_{i-1}),
\label{eq:cbc-enkripsi}
\end{equation}
dan dekripsinya membalik urutan itu,
\begin{equation}
P_i = D_K(C_i) \oplus C_{i-1}.
\label{eq:cbc-dekripsi}
\end{equation}
Blok pertama tidak punya cipherteks sebelumnya, sehingga diperlukan sebuah blok
semu $C_0$ yang disebut \textbf{\textit{Initialization Vector}} (IV). IV dapat
diberikan pengguna atau dibangkitkan secara acak oleh program, dan tidak perlu
dirahasiakan, tetapi tidak boleh dipakai ulang dengan kunci yang sama.

Sifat rantai inilah yang menghapus kelemahan ECB. Dua blok plainteks yang
identik akan di-XOR dengan blok cipherteks pendahulu yang berbeda, sehingga
masukannya bagi $E_K$ berbeda dan hasilnya pun berbeda. Blok $A$ yang muncul di
posisi ke-1 dan ke-4 pada contoh ECB menghasilkan dua blok cipherteks yang
berlainan pada CBC, sehingga pola tetap tersembunyi. Kelemahan manipulasi pun
tertutup: membuang satu blok cipherteks membuat blok berikutnya di-XOR dengan
blok yang salah, sehingga plainteks hasil dekripsi menjadi kacau dan penerima
mengetahui ada yang tidak beres.

Harganya adalah penjalaran galat dan hilangnya sebagian paralelisme. Kesalahan
satu bit pada sebuah blok plainteks saat enkripsi menjalar ke blok cipherteks
yang bersangkutan dan ke seluruh blok cipherteks berikutnya, karena setiap blok
memakai hasil blok sebelumnya. Namun sifat itu berkebalikan pada dekripsi:
kesalahan satu bit pada sebuah blok cipherteks hanya mempengaruhi blok plainteks
yang bersangkutan dan satu bit pada blok plainteks berikutnya, sedangkan blok-blok
sesudahnya kembali benar. Karena itu CBC tetap dapat didekripsi secara paralel
meskipun enkripsinya harus berurutan.

\subsection{Cipher Feedback (CFB)}

Mode CBC menuntut data terkumpul penuh satu blok sebelum dapat dienkripsi,
sehingga kurang cocok untuk saluran yang mengirim satuan kecil, misalnya
per-karakter. Mode \textbf{CFB} mengatasi hal itu dengan memperlakukan cipher
blok seperti cipher alir: data dienkripsi dalam unit $s$ bit, dengan $s \le b$
dan $b$ adalah ukuran blok. Bila unit yang dienkripsi adalah satu karakter
8 bit pada blok 64 bit, modenya disebut \textbf{CFB 8-bit}.

Mode CFB memelihara sebuah antrian (\textit{shift register}) sepanjang $b$ bit.
Antrian itu dienkripsi dengan $E_K$, dan $s$ bit paling kiri dari hasilnya
di-XOR dengan plainteks untuk membentuk cipherteks. Selanjutnya antrian digeser
ke kiri sejauh $s$ bit dan $s$ bit cipherteks yang baru saja terbentuk
disisipkan di sisi kanannya:
\begin{equation}
C_j = P_j \oplus \mathrm{MSB}_s(E_K(X_j)), \qquad X_{j+1} = \mathrm{LSB}_{b-s}(X_j) \parallel C_j,
\label{eq:cfb-enkripsi}
\end{equation}
dengan $X_1 = IV$ dan $\parallel$ menyatakan penyambungan (\textit{concatenation}).
Dekripsinya memakai antrian yang sama,
\begin{equation}
P_j = C_j \oplus \mathrm{MSB}_s(E_K(X_j)), \qquad X_{j+1} = \mathrm{LSB}_{b-s}(X_j) \parallel C_j.
\label{eq:cfb-dekripsi}
\end{equation}
Perhatikan bahwa dekripsi hanya memerlukan fungsi enkripsi $E_K$, bukan $D_K$.

Karena umpan baliknya diambil dari cipherteks, CFB mewarisi sifat penjalaran
galat dari CBC: kesalahan satu bit pada satu blok cipherteks merusak blok
plainteks yang bersangkutan dan menggeser isi antrian sehingga satu blok
berikutnya ikut salah. Dalam kasus khusus $s = b$, seluruh antrian digantikan
oleh cipherteks sebelumnya, sehingga rumusnya menyederhana menjadi
$C_i = P_i \oplus E_K(C_{i-1})$ dengan $C_0 = IV$, yaitu padanan cipher alir
yang persis dari CBC.

\subsection{Output Feedback (OFB)}

Mode \textbf{OFB} hampir sama dengan CFB, dengan satu perbedaan penting pada
arah umpan baliknya. Pada CFB, isi antrian baru diambil dari cipherteks; pada
OFB, isi antrian baru diambil dari keluaran $E_K$ itu sendiri, bukan dari
cipherteks:
\begin{equation}
C_j = P_j \oplus E_K(X_j), \qquad X_{j+1} = E_K(X_j).
\label{eq:ofb}
\end{equation}
Konsekuensinya besar. Karena antrian tidak pernah bergantung pada plainteks
maupun cipherteks, aliran kunci yang dibangkitkan $E_K(X_j)$ dapat dihitung
lebih dahulu, sebelum pesan diketahui. Aliran itu berperilaku persis seperti
keystream pada cipher alir: enkripsi menjadi operasi XOR belaka.

Karena umpan baliknya tidak melewati cipherteks, kesalahan satu bit pada blok
plainteks hanya mempengaruhi blok cipherteks yang bersangkutan, dan kesalahan
satu bit pada blok cipherteks hanya merusak satu bit pada plainteks hasil
dekripsi. Galat tidak menjalar sama sekali. Karakter galat semacam ini cocok
untuk transmisi analog yang didigitisasi, seperti suara atau video, yang dapat
menoleransi sesekali bit yang salah tetapi tidak dapat menoleransi penjalaran
galat ke seluruh pesan. Gambar~\ref{fig:mode-cfb-ofb} memperlihatkan perbedaan
arah umpan balik kedua mode itu.

% Perbandingan arah umpan balik mode CFB dan OFB.
% Dirujuk dari section-02.tex dengan % Perbandingan arah umpan balik mode CFB dan OFB.
% Dirujuk dari section-02.tex dengan % Perbandingan arah umpan balik mode CFB dan OFB.
% Dirujuk dari section-02.tex dengan \input{figures/mode-cfb-ofb}.
\begin{figure}[htbp]
    \centering
    \begin{tikzpicture}[
        kotak/.style={draw, rounded corners, align=center, font=\small, inner sep=3pt,
                      minimum height=0.6cm},
        blok/.style={kotak, fill=blue!6},
        enc/.style={kotak, fill=orange!12, minimum width=1.1cm},
        cip/.style={kotak, fill=green!10},
        xor/.style={draw, circle, inner sep=1pt, font=\small, fill=yellow!18,
                    minimum size=0.5cm},
        panah/.style={-{Latex[length=2mm]}, thick},
        balik/.style={-{Latex[length=2mm]}, thick, dashed},
    ]
        % ================= CFB =================
        \node[font=\small\bfseries] at (2.6,1.5) {Cipher Feedback (CFB)};
        \node[blok, minimum width=2.4cm] (r1) at (0,0.5) {Register $X_j$};
        \node[enc] (e1) at (3.2,0.5) {$E_K$};
        \node[xor] (x1) at (5.4,0.5) {$\oplus$};
        \node[blok] (p1) at (5.4,1.6) {$P_j$};
        \node[cip, minimum width=1.0cm] (c1) at (5.4,-0.6) {$C_j$};
        \draw[panah] (r1) -- node[above, font=\scriptsize] {$b$ bit} (e1);
        \draw[panah] (e1) -- node[above, font=\scriptsize] {$\mathrm{MSB}_s$} (x1);
        \draw[panah] (p1) -- (x1);
        \draw[panah] (x1) -- (c1);
        \draw[balik] (c1.south) -- (5.4,-1.3) -- (-1.5,-1.3) -- (-1.5,0.5) -- (r1.west);
        \node[font=\scriptsize, anchor=west] at (-1.4,-1.65)
            {umpan balik diambil dari cipherteks $C_j$};

        % ================= OFB =================
        \node[font=\small\bfseries] at (2.6,-3.0) {Output Feedback (OFB)};
        \node[blok, minimum width=2.4cm] (r2) at (0,-4.0) {Register $X_j$};
        \node[enc] (e2) at (3.2,-4.0) {$E_K$};
        \node[xor] (x2) at (5.4,-4.0) {$\oplus$};
        \node[blok] (p2) at (5.4,-2.9) {$P_j$};
        \node[cip, minimum width=1.0cm] (c2) at (5.4,-5.1) {$C_j$};
        \draw[panah] (r2) -- (e2);
        \draw[panah] (e2) -- (x2);
        \draw[panah] (p2) -- (x2);
        \draw[panah] (x2) -- (c2);
        \draw[balik] (e2.south) -- (3.2,-5.8) -- (-1.5,-5.8) -- (-1.5,-4.0) -- (r2.west);
        \node[font=\scriptsize, anchor=west] at (-1.4,-6.15)
            {umpan balik diambil dari keluaran $E_K$, bukan dari cipherteks};
    \end{tikzpicture}
    \caption{Arah umpan balik mode CFB dan OFB: CFB mengembalikan cipherteks ke register, sedangkan OFB mengembalikan keluaran fungsi enkripsi sehingga aliran kuncinya tidak bergantung pada data}
    \label{fig:mode-cfb-ofb}
\end{figure}
.
\begin{figure}[htbp]
    \centering
    \begin{tikzpicture}[
        kotak/.style={draw, rounded corners, align=center, font=\small, inner sep=3pt,
                      minimum height=0.6cm},
        blok/.style={kotak, fill=blue!6},
        enc/.style={kotak, fill=orange!12, minimum width=1.1cm},
        cip/.style={kotak, fill=green!10},
        xor/.style={draw, circle, inner sep=1pt, font=\small, fill=yellow!18,
                    minimum size=0.5cm},
        panah/.style={-{Latex[length=2mm]}, thick},
        balik/.style={-{Latex[length=2mm]}, thick, dashed},
    ]
        % ================= CFB =================
        \node[font=\small\bfseries] at (2.6,1.5) {Cipher Feedback (CFB)};
        \node[blok, minimum width=2.4cm] (r1) at (0,0.5) {Register $X_j$};
        \node[enc] (e1) at (3.2,0.5) {$E_K$};
        \node[xor] (x1) at (5.4,0.5) {$\oplus$};
        \node[blok] (p1) at (5.4,1.6) {$P_j$};
        \node[cip, minimum width=1.0cm] (c1) at (5.4,-0.6) {$C_j$};
        \draw[panah] (r1) -- node[above, font=\scriptsize] {$b$ bit} (e1);
        \draw[panah] (e1) -- node[above, font=\scriptsize] {$\mathrm{MSB}_s$} (x1);
        \draw[panah] (p1) -- (x1);
        \draw[panah] (x1) -- (c1);
        \draw[balik] (c1.south) -- (5.4,-1.3) -- (-1.5,-1.3) -- (-1.5,0.5) -- (r1.west);
        \node[font=\scriptsize, anchor=west] at (-1.4,-1.65)
            {umpan balik diambil dari cipherteks $C_j$};

        % ================= OFB =================
        \node[font=\small\bfseries] at (2.6,-3.0) {Output Feedback (OFB)};
        \node[blok, minimum width=2.4cm] (r2) at (0,-4.0) {Register $X_j$};
        \node[enc] (e2) at (3.2,-4.0) {$E_K$};
        \node[xor] (x2) at (5.4,-4.0) {$\oplus$};
        \node[blok] (p2) at (5.4,-2.9) {$P_j$};
        \node[cip, minimum width=1.0cm] (c2) at (5.4,-5.1) {$C_j$};
        \draw[panah] (r2) -- (e2);
        \draw[panah] (e2) -- (x2);
        \draw[panah] (p2) -- (x2);
        \draw[panah] (x2) -- (c2);
        \draw[balik] (e2.south) -- (3.2,-5.8) -- (-1.5,-5.8) -- (-1.5,-4.0) -- (r2.west);
        \node[font=\scriptsize, anchor=west] at (-1.4,-6.15)
            {umpan balik diambil dari keluaran $E_K$, bukan dari cipherteks};
    \end{tikzpicture}
    \caption{Arah umpan balik mode CFB dan OFB: CFB mengembalikan cipherteks ke register, sedangkan OFB mengembalikan keluaran fungsi enkripsi sehingga aliran kuncinya tidak bergantung pada data}
    \label{fig:mode-cfb-ofb}
\end{figure}
.
\begin{figure}[htbp]
    \centering
    \begin{tikzpicture}[
        kotak/.style={draw, rounded corners, align=center, font=\small, inner sep=3pt,
                      minimum height=0.6cm},
        blok/.style={kotak, fill=blue!6},
        enc/.style={kotak, fill=orange!12, minimum width=1.1cm},
        cip/.style={kotak, fill=green!10},
        xor/.style={draw, circle, inner sep=1pt, font=\small, fill=yellow!18,
                    minimum size=0.5cm},
        panah/.style={-{Latex[length=2mm]}, thick},
        balik/.style={-{Latex[length=2mm]}, thick, dashed},
    ]
        % ================= CFB =================
        \node[font=\small\bfseries] at (2.6,1.5) {Cipher Feedback (CFB)};
        \node[blok, minimum width=2.4cm] (r1) at (0,0.5) {Register $X_j$};
        \node[enc] (e1) at (3.2,0.5) {$E_K$};
        \node[xor] (x1) at (5.4,0.5) {$\oplus$};
        \node[blok] (p1) at (5.4,1.6) {$P_j$};
        \node[cip, minimum width=1.0cm] (c1) at (5.4,-0.6) {$C_j$};
        \draw[panah] (r1) -- node[above, font=\scriptsize] {$b$ bit} (e1);
        \draw[panah] (e1) -- node[above, font=\scriptsize] {$\mathrm{MSB}_s$} (x1);
        \draw[panah] (p1) -- (x1);
        \draw[panah] (x1) -- (c1);
        \draw[balik] (c1.south) -- (5.4,-1.3) -- (-1.5,-1.3) -- (-1.5,0.5) -- (r1.west);
        \node[font=\scriptsize, anchor=west] at (-1.4,-1.65)
            {umpan balik diambil dari cipherteks $C_j$};

        % ================= OFB =================
        \node[font=\small\bfseries] at (2.6,-3.0) {Output Feedback (OFB)};
        \node[blok, minimum width=2.4cm] (r2) at (0,-4.0) {Register $X_j$};
        \node[enc] (e2) at (3.2,-4.0) {$E_K$};
        \node[xor] (x2) at (5.4,-4.0) {$\oplus$};
        \node[blok] (p2) at (5.4,-2.9) {$P_j$};
        \node[cip, minimum width=1.0cm] (c2) at (5.4,-5.1) {$C_j$};
        \draw[panah] (r2) -- (e2);
        \draw[panah] (e2) -- (x2);
        \draw[panah] (p2) -- (x2);
        \draw[panah] (x2) -- (c2);
        \draw[balik] (e2.south) -- (3.2,-5.8) -- (-1.5,-5.8) -- (-1.5,-4.0) -- (r2.west);
        \node[font=\scriptsize, anchor=west] at (-1.4,-6.15)
            {umpan balik diambil dari keluaran $E_K$, bukan dari cipherteks};
    \end{tikzpicture}
    \caption{Arah umpan balik mode CFB dan OFB: CFB mengembalikan cipherteks ke register, sedangkan OFB mengembalikan keluaran fungsi enkripsi sehingga aliran kuncinya tidak bergantung pada data}
    \label{fig:mode-cfb-ofb}
\end{figure}


\subsection{Counter Mode (CTR)}

Mode \textbf{Counter} tidak merantai blok sama sekali. Sebagai gantinya, sebuah
nilai \textit{counter} berukuran sama dengan blok plainteks dienkripsi, dan
hasilnya di-XOR dengan plainteks:
\begin{equation}
C_j = P_j \oplus E_K(T_j), \qquad T_j = T_0 + j,
\label{eq:ctr-enkripsi}
\end{equation}
dengan $T_j$ adalah nilai counter untuk blok ke-$j$ dan $T_0$ nilai awalnya.
Nilai counter harus berbeda untuk setiap blok yang dienkripsi; pada praktiknya
counter dinaikkan satu langkah setiap blok. Karena $E_K(T_j)$ tidak bergantung
pada plainteks maupun cipherteks, aliran kunci dapat dihitung lebih dahulu,
seperti pada OFB.

Mode CTR mengumpulkan hampir semua kelebihan mode-mode lain. Hanya fungsi
enkripsi $E_K$ yang dibutuhkan, baik untuk enkripsi maupun dekripsi. Karena
tidak ada ketergantungan antar blok, seluruh blok dapat dienkripsi atau
didekripsi secara paralel, dan blok ke-$j$ dapat langsung diproses tanpa
menghitung blok-blok sebelumnya --- kemampuan akses acak yang pada cipher alir
biasa tidak dimiliki. Kesalahan satu bit hanya merusak satu bit, sama seperti
cipher alir. Gambar~\ref{fig:mode-ctr} memperlihatkan susunan blok yang
mandiri tersebut.

% Mode Counter (CTR): counter dienkripsi lalu di-XOR dengan plainteks.
% Dirujuk dari section-02.tex dengan % Mode Counter (CTR): counter dienkripsi lalu di-XOR dengan plainteks.
% Dirujuk dari section-02.tex dengan % Mode Counter (CTR): counter dienkripsi lalu di-XOR dengan plainteks.
% Dirujuk dari section-02.tex dengan \input{figures/mode-ctr}.
\begin{figure}[htbp]
    \centering
    \begin{tikzpicture}[
        kotak/.style={draw, rounded corners, align=center, font=\small, inner sep=3pt,
                      minimum width=1.35cm, minimum height=0.6cm},
        blok/.style={kotak, fill=blue!6},
        enc/.style={kotak, fill=orange!12, minimum width=1.1cm},
        cip/.style={kotak, fill=green!10},
        xor/.style={draw, circle, inner sep=1pt, font=\small, fill=yellow!18,
                    minimum size=0.5cm},
        grup/.style={draw, dashed, gray!60, rounded corners, inner sep=5pt},
        panah/.style={-{Latex[length=2mm]}, thick},
    ]
        \foreach \i/\x/\t in {1/1.9/{ctr}, 2/4.9/{ctr$+1$}, 3/7.9/{ctr$+2$}} {
            \node[blok] (k\i) at (\x,0)    {\t};
            \node[enc]  (e\i) at (\x,-1.2) {$E_K$};
            \node[xor]  (x\i) at (\x,-2.4) {$\oplus$};
            \node[cip]  (c\i) at (\x,-3.7) {$C_{\i}$};
            \node[blok] (p\i) at (\x-2.0,-2.4) {$P_{\i}$};
            \draw[panah] (k\i) -- (e\i);
            \draw[panah] (e\i) -- (x\i);
            \draw[panah] (p\i) -- (x\i);
            \draw[panah] (x\i) -- (c\i);
            % Kotak putus-putus menandai satu blok yang diproses mandiri,
            % supaya jelas P_i hanya menuju XOR pada panelnya sendiri.
            \node[grup, fit=(k\i)(c\i)(p\i)] {};
        }
        \node[font=\scriptsize, align=center] at (4.9,-4.5)
            {tidak ada ketergantungan antar blok: seluruh blok dapat diproses secara paralel};
    \end{tikzpicture}
    \caption{Mode Counter (CTR): nilai counter yang terus bertambah dienkripsi lalu di-XOR dengan plainteks}
    \label{fig:mode-ctr}
\end{figure}
.
\begin{figure}[htbp]
    \centering
    \begin{tikzpicture}[
        kotak/.style={draw, rounded corners, align=center, font=\small, inner sep=3pt,
                      minimum width=1.35cm, minimum height=0.6cm},
        blok/.style={kotak, fill=blue!6},
        enc/.style={kotak, fill=orange!12, minimum width=1.1cm},
        cip/.style={kotak, fill=green!10},
        xor/.style={draw, circle, inner sep=1pt, font=\small, fill=yellow!18,
                    minimum size=0.5cm},
        grup/.style={draw, dashed, gray!60, rounded corners, inner sep=5pt},
        panah/.style={-{Latex[length=2mm]}, thick},
    ]
        \foreach \i/\x/\t in {1/1.9/{ctr}, 2/4.9/{ctr$+1$}, 3/7.9/{ctr$+2$}} {
            \node[blok] (k\i) at (\x,0)    {\t};
            \node[enc]  (e\i) at (\x,-1.2) {$E_K$};
            \node[xor]  (x\i) at (\x,-2.4) {$\oplus$};
            \node[cip]  (c\i) at (\x,-3.7) {$C_{\i}$};
            \node[blok] (p\i) at (\x-2.0,-2.4) {$P_{\i}$};
            \draw[panah] (k\i) -- (e\i);
            \draw[panah] (e\i) -- (x\i);
            \draw[panah] (p\i) -- (x\i);
            \draw[panah] (x\i) -- (c\i);
            % Kotak putus-putus menandai satu blok yang diproses mandiri,
            % supaya jelas P_i hanya menuju XOR pada panelnya sendiri.
            \node[grup, fit=(k\i)(c\i)(p\i)] {};
        }
        \node[font=\scriptsize, align=center] at (4.9,-4.5)
            {tidak ada ketergantungan antar blok: seluruh blok dapat diproses secara paralel};
    \end{tikzpicture}
    \caption{Mode Counter (CTR): nilai counter yang terus bertambah dienkripsi lalu di-XOR dengan plainteks}
    \label{fig:mode-ctr}
\end{figure}
.
\begin{figure}[htbp]
    \centering
    \begin{tikzpicture}[
        kotak/.style={draw, rounded corners, align=center, font=\small, inner sep=3pt,
                      minimum width=1.35cm, minimum height=0.6cm},
        blok/.style={kotak, fill=blue!6},
        enc/.style={kotak, fill=orange!12, minimum width=1.1cm},
        cip/.style={kotak, fill=green!10},
        xor/.style={draw, circle, inner sep=1pt, font=\small, fill=yellow!18,
                    minimum size=0.5cm},
        grup/.style={draw, dashed, gray!60, rounded corners, inner sep=5pt},
        panah/.style={-{Latex[length=2mm]}, thick},
    ]
        \foreach \i/\x/\t in {1/1.9/{ctr}, 2/4.9/{ctr$+1$}, 3/7.9/{ctr$+2$}} {
            \node[blok] (k\i) at (\x,0)    {\t};
            \node[enc]  (e\i) at (\x,-1.2) {$E_K$};
            \node[xor]  (x\i) at (\x,-2.4) {$\oplus$};
            \node[cip]  (c\i) at (\x,-3.7) {$C_{\i}$};
            \node[blok] (p\i) at (\x-2.0,-2.4) {$P_{\i}$};
            \draw[panah] (k\i) -- (e\i);
            \draw[panah] (e\i) -- (x\i);
            \draw[panah] (p\i) -- (x\i);
            \draw[panah] (x\i) -- (c\i);
            % Kotak putus-putus menandai satu blok yang diproses mandiri,
            % supaya jelas P_i hanya menuju XOR pada panelnya sendiri.
            \node[grup, fit=(k\i)(c\i)(p\i)] {};
        }
        \node[font=\scriptsize, align=center] at (4.9,-4.5)
            {tidak ada ketergantungan antar blok: seluruh blok dapat diproses secara paralel};
    \end{tikzpicture}
    \caption{Mode Counter (CTR): nilai counter yang terus bertambah dienkripsi lalu di-XOR dengan plainteks}
    \label{fig:mode-ctr}
\end{figure}


Satu syarat mutlak berlaku pada CTR: nilai counter tidak boleh berulang. Bila
counter yang sama dipakai dua kali, dua blok memakai aliran kunci yang identik,
dan XOR kedua cipherteksnya akan menghapus kunci sehingga menyisakan XOR kedua
plainteks --- persis serangan \textit{two-time pad} yang dibahas pada
Bab~\ref{chap:stream-cipher}. Karena itu nilai awal counter harus tidak dapat
diprediksi dan harus berubah untuk setiap pesan.

\subsection{Perbandingan Kelima Mode}

Perbedaan kelima mode operasi itu diringkas pada Tabel~\ref{tab:perbandingan-mode}.
Satu hal yang perlu digarisbawahi: mode operasi bukan soal mana yang paling
kuat, melainkan soal mana yang sesuai dengan pola akses data. ECB cocok untuk
data acak yang diakses per blok asalkan tidak ada blok berulang; CBC cocok untuk
berkas yang dienkripsi utuh; CFB dan OFB menangani data yang lebih kecil daripada
satu blok; dan CTR menawarkan paralelisme serta akses acak seluas ECB tanpa
kelemahan polanya. Dalam praktik modern, mode terautentikasi seperti GCM
dibangun di atas prinsip CTR dengan penambahan penanda keutuhan, sebagaimana
dibahas pada Bab~\ref{chap:hash-mac}.

\begin{table}[htbp]
\centering
\small
\begin{tabularx}{\linewidth}{@{}>{\raggedright\arraybackslash}p{1.4cm} >{\raggedright\arraybackslash}p{3.1cm} >{\raggedright\arraybackslash}p{3.7cm} >{\raggedright\arraybackslash}X@{}}
\hline
\textbf{Mode} & \textbf{Ketergantungan} & \textbf{Propagasi galat} & \textbf{Kegunaan khas} \\
\hline
ECB & Tidak ada; tiap blok mandiri &
Galat terbatas pada satu blok &
Data acak yang diakses per blok; pola berulang bocor \\
\hline
CBC & Berantai melalui cipherteks sebelumnya &
Enkripsi menjalar ke semua blok berikutnya; dekripsi hanya 2 blok &
Enkripsi berkas utuh; memerlukan IV \\
\hline
CFB & Berantai melalui cipherteks sebelumnya &
Menjalar sekitar satu blok berikutnya &
Data yang lebih kecil daripada satu blok \\
\hline
OFB & Tidak bergantung cipherteks &
Tidak menjalar; hanya bit yang bersangkutan &
Transmisi suara atau video terdigitalisasi \\
\hline
CTR & Tidak ada; tiap blok mandiri &
Galat terbatas pada bit yang bersangkutan &
Paralel penuh dan akses acak; counter wajib unik \\
\hline
\end{tabularx}
\caption{Perbandingan lima mode operasi cipher blok}
\label{tab:perbandingan-mode}
\sumber{Data dari Munir, \textit{IF4020 Kriptografi}, ITB; dilengkapi dari \citet{stallings2017} dan \citet{menezes1996}.}
\end{table}
